\ifdefined\iscombined
	\lecturetitle{Stability, Impedance, and Kirchhoff's Laws}
\else
\documentclass{article}
\usepackage{xparse}
\usepackage{changepage}
\usepackage{listings}
\usepackage{amsthm}
\usepackage{comment}
\usepackage{amsmath}
\usepackage{braket}
\usepackage{amsfonts}
\usepackage{sectsty}
\usepackage{hyperref}
\usepackage{fancyhdr}
\usepackage[top=1in,bottom=1in,left=1in,right=1in]{geometry}
\usepackage{float}
\usepackage{graphicx}
\usepackage{mathtools}
\usepackage{tikz}
\usetikzlibrary{arrows}

\date{
	\vspace{-.75em}
	{\large \today}
}

\title{
	\vspace*{-1.5em}
	{\Huge Lecture 5} \\
	\vspace{-1em}
	\rule{\textwidth/4}{.01em} \\
	\vspace{-.25em}
	{\Large Stability, Impedance, and Kirchhoff's Laws}
	\vspace{-.75em}
}

\author{
	{\large Matthew Farah}
}

\hypersetup{
	colorlinks=true,
	linkcolor=blue,
	filecolor=magenta,
	urlcolor=blue,
	citecolor=blue,
	anchorcolor=blue
}

\fancyhead{}
\fancyhead[L]{\today}
\fancyhead[R]{Lecture 5 - Stability, Impedance, and Kirchhoff's Laws}

\setlength{\parindent}{0cm}

\newcounter{example}
\renewcommand{\theexample}{\arabic{example}}

\newenvironment{example}[1][]{
	\refstepcounter{example}
	\begin{adjustwidth}{1.5em}{}
	\noindent\textbf{\ifx&#1&Example \theexample:\else Example \theexample \ (#1):\fi}
	\ignorespaces
}{
	\vspace{.5em}
	\end{adjustwidth}
}

\newenvironment{solution}{
	\vspace{.5em}
	\par
	\begin{adjustwidth}{1.5em}{}
	\textbf{{Solution:}}
	\par
	\vspace{.5em}
}{
	\vspace{1em}
	\end{adjustwidth}
}

\newcounter{subexample}
\renewcommand{\thesubexample}{\alph{subexample}}

\newenvironment{subexample}{
	\setcounter{step}{0}
	\refstepcounter{subexample}
	\vspace{1em}\par\noindent
	\begin{adjustwidth}{1.5em}{}
	\noindent\textbf{(\thesubexample)}
	\ignorespaces
}{
	\par
	\vspace{1em}
	\end{adjustwidth}
}

\renewenvironment{proof}{
	\vspace{.5em}\par\noindent
	\begin{adjustwidth}{1.5em}{}
	\emph{Proof:}\par\vspace{1em}
	\setlength{\parindent}{0cm}
}{
	\end{adjustwidth}
}

\newtheorem{theorem}{Theorem}

\renewenvironment{theorem}[1][]{
	\refstepcounter{theorem}
	\vspace{1em}\par\noindent
	\begin{adjustwidth}{1.5em}{}
	\noindent\textbf{\ifx&#1&Theorem \thetheorem:\else Theorem \thetheorem \ (#1):\fi}
	\ignorespaces
}{
	\vspace{.5em}
	\end{adjustwidth}
}

\newtheorem{lemma}{Lemma}

\renewenvironment{lemma}[1][]{
	\refstepcounter{lemma}
	\vspace{1em}\par\noindent
	\begin{adjustwidth}{1.5em}{}
	\noindent\textbf{\ifx&#1&Lemma \thelemma:\else Lemma \thelemma \ (#1):\fi}
	\ignorespaces
}{
	\vspace{.5em}
	\end{adjustwidth}
}

\makeatother
\makeatletter

\newcounter{step}
\renewcommand{\thestep}{\Roman{step}}

\newenvironment{step}{
	\refstepcounter{step}
	\vspace{1em}\par\noindent
	\ignorespaces\noindent\textbf{(\thestep)}
	\ignorespaces
}{
	\par
}

\sectionfont{\fontsize{12}{12}\bfseries}
\subsectionfont{\fontsize{10}{12}\normalfont}
\subsubsectionfont{\fontsize{10}{12}\normalfont}

\begin{document}

\maketitle
\pagestyle{fancy}

\fi

\begin{itemize}
	\item For a system to be stable:
	\begin{enumerate}
		\item It cannot have any poles with positive real components.
		\item It can only have at most one pole at the origin.
	\end{enumerate}
	%FIX: Needs fixing
	\item There exist multiple different definitions of stability:
	\begin{enumerate}
		\item Bounded-Input Bounded-Output (BIBO) stability.
		\item \emph{Note to self: Elaborate}.
	\end{enumerate}
	\item The act of making the real component of the poles of a transfer function more positive, while still remaining negative:
	\begin{enumerate}
		\item Keeps the system stable.
		\item Increases the system's response time.
	\end{enumerate}
	\item Ohm's law states that $v(t) = i(t)R$.
	\begin{itemize}
		\item In the Laplace domain, Ohm's law states that $V(s) = I(s)R$.
	\end{itemize}
	\item The concept of impedance (denoted $Z(s)$) generalizes the concept of resistance to circuit elements other than resistors to make modelling electrical systems easier.
	\begin{itemize}
		\item Impedance is given by:
\begin{align*}
	Z(s) &= \frac{V(s)}{I(s)} = R \\
\end{align*}
	\end{itemize}
	\item Capacitors, inductors, and resistors are all said to have an impedance.
	\item Admittance is defined as the reciprocal of impedance.
	\begin{itemize}
		\item Note: Not used much if at all in this course.
	\end{itemize}
	\item For a capacitor of capacitance $C$, an inductor of inductance $L$, and a resistor of resistance $R$:
	\begin{itemize}
		\item The capacitor's impedance is $\frac{1}{Cs}$.
		\item The inductor's impedance is $Ls$.
		\item The resistor's impedance is $R$.
	\end{itemize}
	\item The impedance of circuit elements in series with impedances $Z_1, \ldots, Z_n$ is given by:
\begin{align*}
	Z_s(s) &= Z_1(s) + \ldots + Z_n(s) \\
\end{align*}
	\item The impedance of circuit elements in parallel with impedances $Z_1, \ldots, Z_n$ is given by:
\begin{align*}
	\frac{1}{Z_p(s)} &= \frac{1}{Z_1(s)} + \ldots + \frac{1}{Z_n(s)} \\
\end{align*}

	\item Kirchhoff's current law states that the sum of all currents entering and leaving a node is zero.
	\begin{itemize}
		\item Current flowing into a node is said to be positive.
		\item Current flowing out of a node is said to be negative.
	\end{itemize}
	\item Kirchhoff's voltage law states that the sum of all voltages around a closed path sum to zero.
\end{itemize}

\ifdefined\iscombined
	\newpage
\else
	\end{document}
\fi
